\notechapter{N-20220411}{Invertibility, Orthogonal Motion, and Polar Decomposition}

A matrix may be invertible while compressing one direction so severely that recovering it from noisy data becomes impossible in practice. Gaussian elimination detects whether a direction is destroyed; singular values quantify how close that destruction is. Between these two questions lies the geometry of orthogonal transformations and polar decomposition.\footnote{See MIT 18.700, supplementary notes on Gaussian elimination and the spectral theorem: \url{https://ocw.mit.edu/courses/18-700-linear-algebra-fall-2013/pages/study-materials/}.}

\section{An inverse requires a trivial kernel}

For \(A\in M_n(F)\), invertibility is equivalent to \(\ker A=\{0\}\). Indeed, rank--nullity shows that an injective endomorphism of \(F^n\) must also be surjective. Gaussian elimination detects this by seeking a pivot in every column. For the two-dimensional case,
\[
 \begin{pmatrix}a&b\\c&d\end{pmatrix}^{-1}
 =\frac1{ad-bc}\begin{pmatrix}d&-b\\-c&a\end{pmatrix},
 \quad ad-bc\ne0.
\]
The determinant is the signed volume factor, so a zero determinant means that some nonzero direction has been collapsed.

The inverse can also be found without a closed formula. To solve $AX=I$, use the same row operations on both blocks of the augmented matrix:
\[
 \left[\begin{array}{cc|cc}
  1&4&1&0\\ 2&9&0&1
 \end{array}\right]
 \;\xrightarrow{R_2\gets R_2-2R_1}\;
 \left[\begin{array}{cc|cc}
  1&4&1&0\\ 0&1&-2&1
 \end{array}\right]
 \;\xrightarrow{R_1\gets R_1-4R_2}\;
 \left[\begin{array}{cc|cc}
  1&0&9&-4\\ 0&1&-2&1
 \end{array}\right].
\]
Thus $A=\left(\begin{smallmatrix}1&4\\2&9\end{smallmatrix}\right)$ has
$A^{-1}=\left(\begin{smallmatrix}9&-4\\-2&1\end{smallmatrix}\right)$.
Every row operation is multiplication by an invertible elementary matrix on the left. Writing the entire computation as $E_k\cdots E_1A=I$ explains both why the right block becomes $A^{-1}=E_k\cdots E_1$ and why elimination is an algorithm for constructing the inverse.

The order of composition is worth checking once:
$(AB)^{-1}=B^{-1}A^{-1}$, because reversing a sequence of operations requires undoing the last one first. In contrast, a diagonal matrix $\operatorname{diag}(d_1,\ldots,d_n)$ is inverted entry by entry if every $d_i$ is nonzero.

\section{The adjugate records the missing direction}

The adjugate \(\operatorname{adj}(A)\) is the transpose of the cofactor matrix. Expanding the determinant along a row gives the identity
\[
 A\operatorname{adj}(A)=\operatorname{adj}(A)A=(\det A)I.
\]
When \(\det A\ne0\), this produces \(A^{-1}\). If \(\operatorname{rank}A=n-1\), some \((n-1)\)-minor is nonzero, so \(\operatorname{adj}(A)\ne0\); the identity then shows every column of the adjugate is a kernel vector. In fact, its rank is one. If \(\operatorname{rank}A\le n-2\), all these minors vanish and \(\operatorname{adj}(A)=0\).

The adjugate is algebraic. The \emph{adjoint} \(A^*\) instead comes from an inner product; it is \(A^T\) over \(\mathbb R\) and \(\overline A^{\,T}\) over \(\mathbb C\). Distinguishing the two matters when interpreting geometry.

\section{Rotations as one-parameter subgroups}

An orthogonal matrix \(Q\) satisfies \(Q^TQ=I\) and preserves all Euclidean lengths and angles. The two-dimensional rotations are
\[
 R_\theta=
 \begin{pmatrix}\cos\theta&-\sin\theta\\
 \sin\theta&\cos\theta\end{pmatrix},
 \qquad R_\theta R_\varphi=R_{\theta+\varphi}.
\]
Differentiating at \(\theta=0\) gives
\[
 J=\begin{pmatrix}0&-1\\1&0\end{pmatrix},\qquad J^T=-J.
\]
This is a Lie-algebraic phenomenon that can be checked directly. If \(K^T=-K\), then \(Q(t)=e^{tK}\) obeys
\[
 \frac d{dt}(Q(t)^TQ(t))=Q(t)^T(K^T+K)Q(t)=0.
\]
Since \(Q(0)=I\), every \(Q(t)\) is orthogonal; continuity of determinant shows \(\det Q(t)=1\). Thus skew-symmetric matrices generate smooth paths of rotations.

\begin{figure}[H]
 \centering
 \begin{tikzpicture}[scale=1.45,>=Stealth]
  \draw[->,gray!75] (-.17,0)--(2.1,0) node[right] {$x$};
  \draw[->,gray!75] (0,-.17)--(0,1.95) node[above] {$y$};
  \coordinate (v) at ({1.72*cos(16)},{1.72*sin(16)});
  \coordinate (w) at ({1.72*cos(54)},{1.72*sin(54)});
  \draw[->,very thick,teal!75!black] (0,0)--(v)
       node[below right] {$v$};
  \draw[->,very thick,orange!85!black] (0,0)--(w)
       node[above left] {$R_\theta v$};
  \draw[->,thick] ({.56*cos(16)},{.56*sin(16)})
       arc[start angle=16,end angle=54,radius=.56];
  \node at (.70,.40) {$\theta$};
 \end{tikzpicture}
 \caption{A rotation changes the direction of a vector while preserving its length; the angle between the two rays is $\theta$.}
\end{figure}

\section{The polar factorization of an invertible map}

Let \(A\in GL(n,\mathbb R)\). The matrix \(A^TA\) is positive definite, so by the spectral theorem it has a unique positive-definite square root
\[
 P=(A^TA)^{1/2}.
\]
Define \(Q=AP^{-1}\). Then
\[
 Q^TQ=P^{-1}A^TAP^{-1}=I,
 \qquad \boxed{A=QP}.
\]
The factor \(P\) is a stretching along orthogonal eigendirections; \(Q\) is an isometry, possibly reversing orientation. Because \(P^2=A^TA\), the two factors are uniquely determined.

For the shear \(A=\left(\begin{smallmatrix}1&1\\0&1\end{smallmatrix}\right)\), the decomposition is explicit:
\[
 \begin{aligned}
 A^TA&=\begin{pmatrix}1&1\\1&2\end{pmatrix},\\
 P&=\frac1{\sqrt5}\begin{pmatrix}2&1\\1&3\end{pmatrix},
 &Q&=\frac1{\sqrt5}\begin{pmatrix}2&1\\-1&2\end{pmatrix}.
 \end{aligned}
\]
Indeed, \(P^2=A^TA\), \(P\) has positive eigenvalues, \(Q^TQ=I\), and multiplying \(QP\) recovers the original shear. Here
\(Q\) is a clockwise rotation by \(\arctan(1/2)\): even a shear has a genuine rotational factor when described using the Euclidean metric. Polar decomposition lets us quantify that geometry instead of attaching a vague arrow to a matrix.

\section{Conditioning is more precise than the determinant}

The singular values of \(A\) are the square roots of the eigenvalues of \(A^TA\), ordered
\(\sigma_1\ge\cdots\ge\sigma_n>0\). They give the sharp estimates
\[
 \sigma_n\|x\|_2\le\|Ax\|_2\le\sigma_1\|x\|_2,
 \qquad \|A^{-1}\|_2=\sigma_n^{-1}.
\]
A perturbation in the data of \(Ax=b\) produces
\[
 \|\delta x\|_2=\|A^{-1}\delta b\|_2
 \le\frac{\|\delta b\|_2}{\sigma_n}.
\]
For relative errors the spectral condition number is
\[
 \kappa_2(A)=\frac{\sigma_1}{\sigma_n}.
\]

The determinant can remain harmless even when the condition number diverges. For
\[
 A_\varepsilon=\begin{pmatrix}\varepsilon&0\\0&\varepsilon^{-1}\end{pmatrix},
 \qquad 0<\varepsilon<1,
\]
we have \(\det A_\varepsilon=1\) but \(\kappa_2(A_\varepsilon)=\varepsilon^{-2}\). The map preserves area while nearly destroying information along one axis.

Elimination determines whether an inverse exists, polar decomposition isolates stretching from orthogonal motion, and singular values show how reliable that inverse is. These are progressively finer questions about a single linear transformation.
